\documentclass[../main.tex]{subfiles}
\graphicspath{{\subfix{../assets/}}}

\begin{document}
\subsection{Copper-based metals}
\subsubsection{Elemental Copper} % 120
Copper as an element is second only to silver in conductivity and exhibits excellent ductility and malleability, all properties highly desirable in electrical applications.

\begin{figure}[h!]
	\begin{center}
		\includegraphics[width=0.75\textwidth]{copperwires.jpg}
	\end{center}
	\caption{Some copper wires \parencite{copper1}}
\end{figure}

Copper isn't just found in wires, however. It finds use in various components ranging from electrical terminals to motors, but all of these take advantage of a few major properties of copper:
\begin{itemize}
	\item Electrical \& thermal conductivity
	\item Ductility \& malleability
	\item Chemical resistance
	\item Infinite recyclability
\end{itemize}
\parencite{copper2}

While conductivity is naturally the paramount property to consider for electrical applications, the other properties elevate copper to an invaluable engineering material.
Not only does its high ductility and recyclability lend itself to easy manufacturing and use, but its chemical resistance ensures longevity once installed, withstanding even extreme industrial loads.
\newpage

\subsubsection{Brass} % 110
While brass measures objectively lower in thermal and electrical conductivity when compared to its parent metal, copper, it makes up this deficit through its higher strength.
Similar to copper, brass is also utilised for wires and endpoint terminals, as well as switch connectors.

\begin{figure}[h!]
	\begin{center}
		\includegraphics[width=0.75\textwidth]{brasswire.png}
	\end{center}
	\caption{Bundled brass wire \parencite{brass2npic}}
\end{figure}

A standard nominal composition for brass is as follows:
\begin{itemize}
	\item $\sim70\%\mathrm{Cu}$
	\item $\sim30\%\mathrm{Zn}$
\end{itemize}
\parencite{brassapps}

While other alloying additions may be varied to adjust the final properties, this proves a baseline of expectation; brass is also more inexpensive compared to pure copper, proving to be an important trait in today's economy.
The extra hardness of the 70-30 alloy also enables easy machining of components, allowing for intricate geometries without structurally unsound parts \textit{e.g.} CPU pins.
An image is included below for clarity.

\begin{figure}[h!]
	\begin{center}
		\includegraphics[width=0.75\textwidth]{cpupins.jpg}
	\end{center}
	\caption{Brass CPU pins \parencite{cpupins}}
\end{figure}

\newpage

\subsubsection{Chromium, Beryllium, and Zirconium Copper} % 160
While all 3 of these metals may be alloyed with copper with varied resulting properties, this report will only discuss beryllium copper for the sake of brevity, noting its comparatively excellent strength.

\begin{figure}[h!]
	\begin{center}
		\includegraphics[width=0.75\textwidth]{becu.jpg}
	\end{center}
	\caption{Beryllium copper (BeCu) pieces \parencite{becu1}}
\end{figure}

A typical nominal composition for beryllium copper (C17200) is as follows:
\begin{itemize}
	\item $\sim98\%\mathrm{Cu}$
	\item $\sim2\%\mathrm{Be}$
\end{itemize}
\parencite{becu1}

Boasting strength comparable to alloy steels, beryllium copper is ideal for use in contact components such as switches and relays.
In particular, while it's not utilised for wiring due to its high strength, it still finds use in connectors between components.
Some specific properties and their impacts are as follows:
\begin{itemize}
	\item Comparable strength to steels allows for use in conductive yet load bearing parts
	\item Moderate conductivity reduces, but doesn't eliminate use in electrical parts
\end{itemize}
\parencite{becu2}

While the conductivity of beryllium copper is roughly one-fifth that of pure copper, it retains enough conductivity to justify use.
Contact components specifically require materials capable of withstanding repeated mechanical load, making beryllium copper an ideal compromise between durability and conductivity (efficiency).
\newpage

\subsection{Forming Copper for Electrical Applications}
\subsubsection{Cold Working (drawing, extrusion, \textit{etc.})}
\paragraph{Die Drawing} % 130
The ductile nature of copper and most cuprous alloys lends itself to cold working with relative ease, with drawing and extrusion being prime examples.
Illustrated below is a die for drawing wires.

\begin{figure}[h!]
	\begin{center}
		\begin{subfigure}[t]{0.45\textwidth}
			\begin{center}
				\includegraphics[width=\textwidth]{drawdie.png}
			\end{center}
			\caption{Technical specifications for a wire die \parencite{wiredraw}}
		\end{subfigure}%
		~~
		\begin{subfigure}[t]{0.45\textwidth}
			\begin{center}
				\includegraphics[width=\textwidth]{wiredrawing.png}
			\end{center}
			\caption{Stress diagrams of various drawn wires \parencite{wiredraw}}
		\end{subfigure}
	\end{center}
	\caption{Various wire drawing diagrams}
\end{figure}

The process of die drawing can be summarised as follows:
\begin{enumerate}
	\item The material pre-shaped into a suitable diameter from which to reduce
	\item The wire is sequentially passed through a series of dies with decreasing diameter
	\item As the wire lengthens, the diameter decreases
\end{enumerate}

Notice that the wire is \textit{drawn} through the dies as compared to extrusion, where the material is \textit{pushed} through the dies.
Wires and cables are formed via die drawing for its economic efficiency, while more complex geometries must be formed through machining or extrusion.
However, most, if not all wires used in electrical applications are drawn.
\newpage

\paragraph{Extrusion} % 180
The process of extrusion is both similar and vastly different from die drawing.
While the material is passed through a narrow opening (the extrusion die), it's instead pushed through with hydraulic pressure (typically) while heated above the recrystallisation temperature.

\begin{figure}[h!]
	\begin{subfigure}[t]{0.35\textwidth}
		\includegraphics[width=\textwidth]{extrusionpic.jpg}
		\caption{A diagram of the extrusion process \parencite{extrusion}}
	\end{subfigure}%
  ~~
	\begin{subfigure}[t]{0.6\textwidth}
		\includegraphics[width=\textwidth]{extrusionexamples.png}
		\caption{Examples of extruded parts \parencite{extrusion2}}
	\end{subfigure}
  \caption{Diagram and example of extruded parts}
\end{figure}

While the specific part illustrated in the above figure may not match any specific electrical components, parts like busbars, connectors, and terminals are common candidates for extrusion manufacturing.
Some specific benefits of extrusion include:
\begin{itemize}
	\item Complex, continuous geometry capabilities
	\item Economically efficient
	\item Effective surface finishing
\end{itemize}
If the specification of \textit{electrical applications} is to be taken quite literally, then the aforementioned components highlight the benefits of extrusion, allowing for more complex geometries than die drawing while maintaining part continuity.
However, if this specification is taken more loosely, then examples may also include heat sinks and other components used in electrical applications, albeit not as an electrical conductor.
An image is included below.

\begin{figure}
  \begin{center}
    \includegraphics[width=0.75\textwidth]{heatsink.jpg}
  \end{center}
  \caption{(Extruded) copper heat sink \parencite{heatsink}}
\end{figure}
\newpage

\subsubsection{Casting} % 240
Copper may also be cast for electrical applications, allowing for the production of complex geometries that may be unrealistic to form by other means.
Specifically, casting allows for the production of substantially larger parts compared to other methods without sacrificing desirable mechanical properties utilising varying casting techniques.
For the sake of example, some casting techniques include:
\begin{itemize}
  \item Sand casting 
  \item Continuous casting 
  \item Centrifugal casting 
  \item Various other mould-based methods \textit{e.g.}
    \begin{itemize}
      \item Shell moulding 
      \item Gravity die 
      \item Chill casting
    \end{itemize}
\end{itemize}

This report will discuss continuous and sand casting, specifically.

\paragraph{Continuous casting}
While less common in electrical applications directly, continuous casting is still an essential method in producing material "blanks" for extrusion and die drawing.
The process is rather simple: molten metal is constantly passed through a thermally conductive mould that is continuously cooled to form an uninterrupted part that can then be cut to size or further refined. \parencite{casting1}

\begin{figure}[h!]
  \begin{center}
    \includegraphics[width=0.75\textwidth]{contcast.jpg}
  \end{center}
  \caption{Continuous cast copper \parencite{castingpic2}}
\end{figure}
\newpage

\paragraph{Sand casting}
One of the age-old methods of forming, sand casting is characterised by its well-known procedure and developed toolkit for complex volumes.
These parts are common in high-volume high-intensity applications such as motors and generators, as well as relay connectors where function supersedes form.
Indeed, the issues of poor surface finishing or uneven cooling may be unavoidable, but the flexibility of sand casting is a justifiable compromise. \parencite{casting1}

\begin{figure}[h!]
  \begin{center}
    \includegraphics[width=0.75\textwidth]{castcopper.jpg}
  \end{center}
  \caption{Sand cast bushings \parencite{castingpic}}
\end{figure}

A reference procedure is as follows:
\begin{enumerate}
  \item Reference creation 
  \item Mould fitting \& packing 
  \item Inner void fills 
  \item Add risers and vents 
  \item Melt metal \& pour 
  \item Once cooled: remove from mould and process as necessary
\end{enumerate}
\newpage

\sloppy
\printbibliography[heading=subbibnumbered,keyword={sect2},title={References}]
\fussy
\end{document}
